% !PRINTED 552
\weberpage{552}
Die Werthe von $\cos m\varphi$ und $\sin m\varphi$ denken wir uns kleiner als 1 und so, dass ihre Quadratsumme $= 1$ ist, gegeben.

Wir adjungiren noch $m^{\text{te}}$ Einheitswurzeln, mit denen wir uns im nächsten Abschnitte noch eingehender beschäftigen werden, die aber jedenfalls von Gleichungen abhängen, deren Grad niedriger als $m$ ist. Um aber den Körper reell zu behalten, wollen wir nicht die Einheitswurzeln selbst, sondern
\weq{1}{\sin\frac{2\pi}{m}, \qquad \cos\frac{2\pi}{m}}
adjungiren.

Es möge also der Körper $\Omega$ aus allen rationalen Zahlen und aus den rationalen Functionen von $\cos m\varphi$, $\sin m\varphi$, $\cos\frac{2\pi}{m}$, $\sin\frac{2\pi}{m}$ bestehen.

Die Gleichung $m^{\text{ten}}$ Grades, von der $x = 2\cos\varphi$ abhängt, haben wir in \S.\,137 aufgestellt in der Form
\weq{2}{2\cos m\varphi = A_m(x),}
worin $A_m(x)$ eine ganze Function $m^{\text{ten}}$ Grades von $x$ ist; wir haben aber dort auch noch die Gleichung
\weq{3}{\sin\varphi = \frac{\sin m\varphi}{B_m(x)}}
gefunden, durch die $\sin\varphi$ rational durch $x$ ausgedrückt ist. In diesen Formeln sind die Einheitswurzeln noch nicht enthalten.

Die $m$ Wurzeln von (2) haben nun folgende Bedeutung:
\[
x_0 = 2\cos\varphi, \qquad x_1 = 2\cos\left(\varphi + \frac{2\pi}{m}\right)\dots,
\]
\[
x_{m-1} = 2\cos\left(\varphi + \frac{2(m-1)\pi}{m}\right),
\]
und mit Hülfe von (2) und (3) und unter Adjunction der Grössen (1) kann jede von ihnen als rationale Function einer anderen dargestellt werden, und zwar so:
\[
x_1 = f(x_0), \qquad x_2 = f(x_1)\dots, \qquad x_0 = f(x_{m-1}).
\]

Irgend eine Function der $x_0, x_1 \dots x_{m-1}$ kann also dargestellt werden als rationale Function $F(x_0)$, und bei einer cyklischen Permutation geht
\[
F(x_0) \text{ in } F(x_1), \ F(x_1) \text{ in } F(x_2) \dots F(x_{m-1}) \text{ in } F(x_0)
\]
über. Für eine \gesperrt{cyklische Function} ist also
%CONT>>
